\scanpage{136}
que cette section du torseur (qui est
en fait défini sur $\widehat X$, compactifiée
non singulière naturelle de $X$, \unsure{\}|)} soit
${\unsure{\mathscr C|C|\mathcal C}}_{\widehat X/\mathbb C}$) n'admet ``à l'infini'' $\widehat X-X$
que des pôles simples au plus.

Si $D$ est le diviseur correspondant
\inserted{sur $\widehat X$} à multiplicités toutes $1$, de support $\widehat X-X$,
on aura comme groupe d'indétermination
des conn. proj. rég. à l'infini sur $\widehat X$
\[
 H^0(\widehat X,\underline\omega_{\widehat X}^{\otimes2}(D))
\]
qu'on peut relier ainsi aux invariants classiques
de la courbe complète $\widehat X$ (en écrivant
\deleted{\unsure{maintenant}}\inserted{d'abord} $\deleted{\underline{\omega}^{2}}\ \underline\lambda\deleted{\unsure{\otimes2}}$ au lieu de $\underline{\omega}^{\otimes2}_{\widehat X/\mathbb C}$)
\[
\deleted{0\longrightarrow\underline{\omega}^{2}\longrightarrow
 \underset{\deleted{\unsure{\Theta}}}{\inserted{\underline\lambda(D)}}\longrightarrow
 \underline{\omega}^{2}(D)/\underline{\omega}^{2}\longrightarrow0}
\]
\[
0\longrightarrow\underline\lambda\longrightarrow\underline\lambda(D)
\longrightarrow\underline\lambda(D)/\underline\lambda\longrightarrow0
\]
déduite par $\otimes\,\underline\lambda$ de

\diagram{\begin{tikzcd}[ampersand replacement=\&,column sep=small,row sep=tiny]
0\arrow[r]\&\underline O_{\widehat X}\arrow[r]\&
\underline O_{\widehat X}(D)\arrow[r]\&
\underline O_{\widehat X}(D)/\underline O_{\widehat X}\arrow[r]\&0\\
\&\&\underbrace{\mathscr J_D^{-1}}_{\text{Idéal de }D}\arrow[u,equal]\&
(\mathscr J_D/\mathscr J_D^2)^{-1}\arrow[u,"\unsure{\wr|\simeq}" description]\\
\&\&\&\deleted{\unsure{(\underline\omega|D)^{-}}}\arrow[u,"\unsure{\wr|\simeq}" description]\\[-1ex]
\&\&\&(\underline\omega^{-1}|D)
\end{tikzcd}}

d'où
\[
\underline\lambda(D)/\underline\lambda\simeq(\underline\lambda\,\underline\omega^{-1})|D
\]
et la suite exacte

\diagram{\begin{tikzcd}[ampersand replacement=\&,column sep=tiny,row sep=small]
0\arrow[r]\&H^0(\widehat X,\underline\lambda)\arrow[r]\&
H^0(\widehat X,\underline\lambda(D))\arrow[r]\&
H^0(D,\underline\lambda\,\underline\omega^{-1}|D)\arrow[d]\\
\&\&\&H^1(\widehat X,\underline\lambda).
\end{tikzcd}}

(ici $\underline\lambda=\underline\omega_{\widehat X/\mathbb C}^{\otimes2}\simeq\underline\omega^{2}$

\inserted{On suppose $X$ connexe,} \textcircled{a} Si $\widehat X$ est de genre $g\geq2$, $H^1(\widehat X,\underline\lambda)=0$
(c'est le dual de $H^0(\widehat X,\underline\omega^{-1})=0$) donc on
trouve

\diagram{\begin{tikzcd}[ampersand replacement=\&,column sep=tiny,row sep=small]
0\arrow[r]\&H^0(\widehat X,\underline\omega^2)\arrow[r]\&
H^0(X,\underline\omega^2)^{\text{rég}}\arrow[r]\&
H^0(D,\underline\omega|_{\deleted{\unsure{X}}D})\arrow[r]\&0\\
\&\unsure{:=|=}\Omega^1_{M_g/\mathbb C}\,(s)\arrow[u,"\unsure{\wr|\simeq}" description]\&
\Omega^1_{M_{g,\nu}/\mathbb C}\,(s)\arrow[u,"\unsure{\wr|\simeq}" description,"{\text{can. ?}}"']\&
\text{rang }\nu=\operatorname{card}\deleted{\unsure{D/}}\inserted{\operatorname{supp}D}
\end{tikzcd}}

\margin{rang \deleted{\unsure{g}} $3g-3$}
$M_g$ variété modulaire des courbes de genre $g$
\inserted{lisses complètes} avec $M_{g,\nu}$ variété modulaire \unsure{pour}
pointées à $\nu$ trous
% archive sheet number: 135
